\newpage

% *************************
% **     PROBLEMA        **
% *************************

\section{Problema de Investigación}
\subsection{Descripción del problema}

A nivel global, la educación en lenguas indígenas y minoritarias enfrenta una brecha tecnológica persistente. Más de la mitad de los niños en edad escolar no aprende en su lengua materna durante los primeros años de escolaridad, lo que afecta negativamente su comprensión lectora, su identidad cultural y sus trayectorias educativas a largo plazo (UNESCO, 2022). El acceso a recursos pedagógicos culturalmente pertinentes —adaptados a la lengua y al contexto del estudiante— representa una necesidad urgente y aún no resuelta, especialmente en regiones rurales con limitaciones de infraestructura digital.

En América Latina, esta problemática adquiere una dimensión étnica y lingüística particular. Más de 45 millones de personas pertenecen a pueblos indígenas distribuidos en 23 países, con lenguas que en su mayoría carecen de recursos digitales estructurados para la enseñanza formal. Los estados de la región han avanzado normativamente hacia la Educación Intercultural Bilingüe (EIB) —el modelo pedagógico que busca garantizar el derecho de los estudiantes de pueblos originarios a aprender en su lengua materna junto al español, en diálogo con su cultura y con otras culturas—, pero la brecha entre el mandato normativo y la disponibilidad real de herramientas tecnológicas pertinentes permanece amplia.

En el Perú, esta tensión se manifiesta de manera crítica. El sistema EIB atiende a más de 1,2 millones de estudiantes en aproximadamente 26\,400 instituciones educativas \citep{minedu2018}, concentradas principalmente en el sur andino y la Amazonía. Según datos del INEI, al primer trimestre de 2025 solo el 20,5\% de los hogares del área rural cuenta con acceso a internet, frente al promedio nacional que supera el 58\%. Esta brecha digital se profundiza especialmente en los departamentos del sur andino —Cusco, Apurímac, Ayacucho, Puno y Huancavelica—, precisamente donde se concentra la mayor parte del alumnado EIB, que carece casi por completo de herramientas tecnológicas adaptadas a su realidad lingüística y territorial.

Dentro del sur andino peruano, la lengua indígena predominante es el Quechua Collao, variedad hablada principalmente en Cusco, Puno, Arequipa y Moquegua. Esta variedad posee particularidades lingüísticas de alta relevancia para el diseño de recursos tecnológicos: su morfología aglutinante permite construir palabras de gran complejidad semántica mediante sufijación encadenada, y su fonología incluye tres series de consonantes oclusivas —simples, aspiradas y glotizadas— que no tienen equivalente en el español. Estas características la convierten en una lengua especialmente compleja para los sistemas de procesamiento automático del lenguaje. Adicionalmente, los estudiantes de estas comunidades frecuentemente alternan entre el quechua y el español en una misma interacción —fenómeno denominado translengüaje o alternancia de códigos—, lo que exige que cualquier herramienta de apoyo educativo sea capaz de operar en ambas lenguas y de identificar el idioma predominante de cada consulta.

Frente a este escenario, los avances recientes en inteligencia artificial generativa han abierto posibilidades para el desarrollo de asistentes educativos capaces de responder preguntas, explicar contenidos y apoyar procesos de aprendizaje. Sin embargo, la gran mayoría de estas soluciones presentan tres limitaciones estructurales que las hacen inaplicables en el contexto EIB rural del sur andino: primero, dependen de conectividad estable a internet para acceder a modelos alojados en servidores remotos; segundo, están diseñadas con datos de lenguas de alta disponibilidad, con escasa o nula representación del Quechua Collao; tercero, no están adaptadas a contenidos curriculares ni a la realidad pedagógica de las aulas EIB.

La literatura técnica muestra avances específicos pero fragmentados en este campo. \citet{SullaTorres2024} desarrollaron un modelo de traducción automática neuronal optimizado con BERT para la dirección quechua-español, reportando un BLEU de 20,13 y destacando las dificultades de la morfología aglutinante del quechua para los modelos estándar. \citet{Orcotoma2025} contribuyeron con la creación de un corpus paralelo Quechua Collao-español de 44\,263 palabras usando reconocimiento óptico de caracteres sobre diccionarios de la Academia Mayor de la Lengua Quechua, sentando precedentes para la recopilación de datos en esta variedad específica. \citet{Duran2026} propusieron un método semi-automático de construcción de diccionarios quechua-español utilizando formalismos lingüísticos y transductores de estados finitos, ampliando los recursos léxicos disponibles. Sin embargo, estos trabajos se concentran principalmente en subtareas lingüísticas y no en la construcción de un asistente educativo offline alineado con materiales curriculares y necesidades docentes.

Del mismo modo, los estudios sobre despliegue offline de modelos pequeños demuestran viabilidad técnica. \citet{Walusimbi2026} proponen Arapai, una arquitectura de chatbot educativo offline-first que valida que modelos cuantizados a 4 bits pueden funcionar en dispositivos locales con latencia de 1 a 3 segundos. \citet{Hevia2025} evaluaron un sistema de recuperación aumentada por generación combinado con un modelo pequeño para educación offline, encontrando que la calidad del fragmentado de documentos es crítica para el rendimiento del sistema. No obstante, estos estudios se centran en inglés y no abordan los retos de lenguas aglutinantes de bajos recursos.

El problema se agrava cuando los estudiantes formulan consultas en español, Quechua Collao o mezclas bilingües —fenómeno documentado en aulas EIB del sur andino como alternancia de códigos o translengüaje—. En estos casos, un sistema educativo debe identificar el idioma predominante, reconocer posibles segmentos alternados y recuperar materiales adecuados al contexto. Si el sistema falla en esta etapa, puede responder en una lengua no pertinente, perder información cultural o generar explicaciones desconectadas del currículo local. El reto no es únicamente técnico, sino también pedagógico y lingüístico.

En consecuencia, existe una brecha de investigación y desarrollo: no se identifica una solución integrada de asistente educativo bilingüe español-Quechua Collao que funcione offline, use modelos pequeños de lenguaje, incorpore detección automática del idioma y recupere contenido curricular local para apoyar procesos de aprendizaje en zonas rurales de EIB. Esta brecha justifica el diseño de una propuesta que articule las capas lingüística, pedagógica y tecnológica en un único artefacto evaluable.

\subsection{Identificación del problema}
El problema central de la investigación es la ausencia de un asistente educativo bilingüe español-Quechua Collao, ejecutable sin conexión a internet, que integre detección automática del idioma, recuperación local de contenidos curriculares y generación de respuestas pedagógicas pertinentes para estudiantes y docentes de contextos rurales de Educación Intercultural Bilingüe del sur andino peruano.

\subsection{Formulación del problema}
\subsubsection{Problema general.}
¿Cómo diseñar e implementar un asistente educativo bilingüe español-Quechua Collao basado en un Small Language Model que funcione offline en dispositivos de bajo o mediano costo para apoyar procesos de aprendizaje en contextos rurales de Educación Intercultural Bilingüe del sur andino peruano?

\subsubsection{Sub problemas.}
\begin{itemize}
	\item ¿Cuáles son los requerimientos lingüísticos, pedagógicos y tecnológicos que debe satisfacer un asistente educativo bilingüe offline para contextos rurales de EIB del sur andino peruano?
	\item ¿Cómo diseñar una arquitectura de software que integre detección automática del idioma, recuperación local de contenidos curriculares y generación de respuestas pedagógicas en español y Quechua Collao?
	\item ¿Cómo implementar y configurar un Small Language Model cuantizado para operar sin conexión a internet en dispositivos de bajo o mediano costo con pertinencia lingüística para el Quechua Collao?
	\item ¿Cómo evaluar la calidad lingüística, la utilidad educativa, la adecuación intercultural y el rendimiento técnico offline del asistente propuesto?
\end{itemize}


% *************************
% **     JUSTIFICACION   **
% *************************
\section{Justificación}
La investigación se justifica porque las condiciones descritas en el problema —escasa conectividad rural, ausencia de herramientas tecnológicas en Quechua Collao y más de 1,2 millones de estudiantes EIB desatendidos tecnológicamente \citep{minedu2018}— no encontrarán solución con las herramientas de IA generativa existentes, que dependen de internet y de lenguas de alta disponibilidad. El proyecto propone una respuesta específica a esa brecha desde tres perspectivas complementarias:

\begin{itemize}
	\item \textbf{Justificación Teórica:} El estudio aporta a un campo todavía fragmentado. Mientras existen investigaciones sobre traducción automática neural para quechua \citep{SullaTorres2024, Orcotoma2025} y se sabe que los Grandes Modelos de Lenguaje presentan deficiencias sistemáticas en lenguas de bajos recursos \citep{Court2024}, estas líneas no se han integrado con un enfoque educativo bilingüe. La investigación propone articular avances en inferencia local de modelos pequeños \citep{Xu2024} con detección de idioma y generación aumentada por recuperación (RAG).

	\item \textbf{Justificación Práctica:} El proyecto propone una alternativa directa para apoyar a docentes y estudiantes en entornos de baja conectividad sin necesidad de reemplazar al docente. Su propósito es ofrecer apoyo localizado para responder consultas, explicar conceptos y recuperar información curricular manteniendo el Quechua Collao como lengua válida de aprendizaje. Estudios recientes, como la arquitectura Arapai \citep{Walusimbi2026}, confirman que el despliegue de chatbots offline con modelos pequeños es viable para entornos con limitaciones de infraestructura similares al sur andino peruano.

	\item \textbf{Justificación Metodológica:} La propuesta permite organizar un proceso de diseño y evaluación que combina herramientas de procesamiento de lenguaje natural con criterios pedagógicos. La integración de sistemas RAG \citep{Li2025} y criterios de tutoría inteligente \citep{Levonian2025} permite evaluar el sistema no solo por métricas lingüísticas algorítmicas (como BLEU o chrF++), sino también por criterios de utilidad educativa, pertinencia cultural, y estabilidad \textit{offline} en hardware de bajo costo.
\end{itemize}

La viabilidad del estudio se sostiene en la disponibilidad creciente de modelos compactos y herramientas de inferencia local que pueden ejecutarse mediante técnicas de cuantización \citep{Xu2024}.

% *************************
% **     ALCANCES        **
% *************************
\section{Alcances y limitaciones}
\subsection{Alcances}
El presente estudio abarca el diseño, implementación y evaluación de un asistente educativo bilingüe español-Quechua Collao basado en un Small Language Model. El artefacto principal es un prototipo de sistema informático que no desarrolla un modelo de lenguaje desde cero ni realiza entrenamiento del modelo base, sino que selecciona, cuantiza y configura un SLM pre-entrenado e integra un pipeline de recuperación aumentada por generación con contenidos curriculares EIB locales. La investigación se centra en la asistencia textual para estudiantes y docentes de contextos rurales de EIB, considerando entradas en español, Quechua Collao o mezcla bilingüe. Se definen los siguientes alcances específicos:

\begin{itemize}
	\item \textbf{Alcance lingüístico:} Se delimita de forma exclusiva al Quechua Collao, de acuerdo con la ubicación geográfica del estudio y la disponibilidad de corpus paralelos como los construidos por \citet{Orcotoma2025}. Esta variedad presenta rasgos fonológicos diferenciadores (consonantes oclusivas glotizadas y aspiradas) que deben reflejarse cuidadosamente en los recursos del sistema. Esta delimitación evita prometer cobertura de otras variedades dialectales, permitiendo construir un prototipo culturalmente pertinente.

	\item \textbf{Alcance pedagógico:} El asistente está diseñado para apoyar consultas de nivel primaria (3.° a 6.° grado) y secundaria básica, en las áreas curriculares de Comunicación y Ciencia y Tecnología, conforme al Currículo Nacional (CN-2016) y los materiales EIB del MINEDU. Las respuestas seguirán el principio de andamiaje cognitivo, priorizando explicaciones graduales sobre respuestas directas, y el corpus pedagógico se construirá con libros de texto EIB, el Currículo Nacional y glosarios terminológicos del MINEDU para Quechua Collao.

	\item \textbf{Alcance tecnológico:} Comprende una arquitectura offline compuesta por detección automática del idioma, enrutamiento lingüístico, recuperación aumentada por generación con recursos locales y un Small Language Model pre-entrenado con técnicas de cuantización a 4 bits. El aporte tecnológico radica en la integración, configuración y evaluación del \textit{pipeline} completo.

	\item \textbf{Alcance disciplinar:} Se ubica en la convergencia de la ingeniería informática, el procesamiento de lenguaje natural, la tecnología educativa y la Educación Intercultural Bilingüe. Se orienta como una investigación aplicada de alcance tecnológico-experimental, que implementa y evalúa un prototipo funcional bajo condiciones controladas.
\end{itemize}

\subsection{Limitaciones}
Se reconocen las siguientes limitaciones que condicionan la ejecución del proyecto:

\begin{itemize}
	\item \textbf{Disponibilidad de corpus (Limitación de datos):} Al tratarse de una lengua de bajos recursos digitales, los datos existentes pueden ser reducidos o heterogéneos, lo cual afecta la recuperación y comprensión de los modelos \citep{Court2024}. Para mitigar esto, se complementará el corpus con glosarios, materiales EIB y datos sintéticos validados por hablantes nativos.

	\item \textbf{Restricción dialectal:} El prototipo construido está estrictamente orientado a los rasgos fonológicos y léxicos del Quechua Collao. El sistema no está diseñado para generalizarse o atender consultas en otras variedades quechuas.

	\item \textbf{Restricciones de hardware (Tecnológicas):} La ejecución offline de un SLM depende estrictamente de la memoria disponible (RAM), el \textit{runtime} y el consumo energético del dispositivo local \citep{Xu2024}. Esto restringe el tamaño del modelo utilizable y la complejidad de las consultas que puede atender con una latencia aceptable.

	\item \textbf{Evaluación pedagógica multidimensional:} La utilidad educativa no puede medirse únicamente con métricas automáticas. Las interacciones generativas requieren evaluación de relevancia y pertinencia pedagógica \citep{Levonian2025}, lo que demanda la revisión obligatoria por parte de docentes o especialistas EIB.

	\item \textbf{Acceso territorial para validación:} La evaluación con usuarios reales requiere coordinación institucional con Direcciones Regionales (DRE) y Unidades de Gestión Educativa Local (UGEL), lo cual está condicionado por factores logísticos, distancias geográficas y autorizaciones escolares.

	\item \textbf{Alcance exclusivamente textual:} El prototipo no incorpora reconocimiento ni síntesis de voz. Los desafíos del habla bilingüe se consideran fuera del alcance actual y como posibles líneas de ampliación futura.

	\item \textbf{Dependencia de especialistas institucionales:} La validación final del asistente depende de la disponibilidad y experiencia de especialistas en Quechua Collao vinculados al ámbito institucional. Sin esta validación, el prototipo carecería de la fundamentación cultural necesaria.
\end{itemize}

% *************************
% **     OBJETIVOS       **
% *************************

\section{Objetivos}
\subsection{Objetivo General}
Diseñar, implementar y evaluar un asistente educativo bilingüe español-Quechua Collao basado en un Small Language Model, con funcionamiento offline, detección automática del idioma y recuperación local de contenidos curriculares, para apoyar procesos de aprendizaje en contextos rurales de Educación Intercultural Bilingüe del sur andino peruano.

\subsection{Objetivos Específicos}
\begin{itemize}
	\item Determinar los requerimientos funcionales lingüísticos, pedagógicos y tecnológicos que debe satisfacer un asistente educativo bilingüe offline para contextos rurales de EIB del sur andino peruano, mediante revisión de literatura especializada, análisis de documentos normativos del MINEDU y consulta con especialistas en Quechua Collao.
	\item Diseñar una arquitectura de software offline que integre detección automática del idioma de entrada, enrutamiento lingüístico, recuperación aumentada por generación con base local de contenidos EIB, y generación de respuestas pedagógicas contextualizadas mediante un Small Language Model cuantizado.
	\item Implementar un prototipo funcional del asistente educativo bilingüe, seleccionando, cuantizando y configurando el Small Language Model base e integrando la base local de contenidos curriculares en español y Quechua Collao.
	\item Evaluar el desempeño técnico del prototipo considerando latencia, consumo de memoria, tokens por segundo, estabilidad offline y capacidad de respuesta en dispositivos de bajo o mediano costo.
	\item Evaluar la pertinencia lingüística, utilidad educativa y adecuación intercultural de las respuestas generadas por el asistente, con participación de docentes y especialistas EIB, mediante métricas automáticas de calidad lingüística y rúbricas de valoración pedagógica.
\end{itemize}

\subsection{Alineación con los Objetivos de Desarrollo Sostenible}
La investigación se alinea con los siguientes Objetivos de Desarrollo Sostenible (ODS) de la Agenda 2030 de las Naciones Unidas:

\begin{itemize}
	\item \textbf{ODS 4 — Educación de calidad.} El estudio contribuye a garantizar una educación inclusiva, equitativa y de calidad al proponer una herramienta de apoyo pedagógico en lengua materna para estudiantes de pueblos originarios en zonas rurales. El asistente educativo bilingüe busca reducir las barreras lingüísticas y tecnológicas que limitan el acceso a recursos de aprendizaje en contextos de Educación Intercultural Bilingüe.

	\item \textbf{ODS 9 — Industria, innovación e infraestructura.} La propuesta promueve la innovación tecnológica al diseñar una solución basada en modelos pequeños de lenguaje con funcionamiento offline, adaptada a condiciones de infraestructura limitada. El despliegue de un SLM cuantizado en dispositivos de bajo costo representa una contribución a la infraestructura tecnológica educativa de las comunidades rurales del sur andino.

	\item \textbf{ODS 10 — Reducción de las desigualdades.} El proyecto aborda directamente la brecha digital y lingüística que afecta a las comunidades quechuahablantes. Al ofrecer un asistente que funciona sin conexión a internet y que opera en Quechua Collao, se busca reducir la desigualdad en el acceso a herramientas educativas tecnológicas entre zonas urbanas y rurales.
\end{itemize}
